\hypertarget{appendix-46da295c6787}{}\label{referencestart:appendix-46da295c6787}
\section{Covectors linear functionals and normal transforms}\label{appendix-46da295c6787-covectors-linear-functionals-and-normal-transforms}

\textbf{Curriculum depth:} Developed · \textbf{First seen:} Phase6 · \textbf{Latest development:} Phase7

\begin{allmightroledefinition}

\subsection{Definition and mechanism}\label{appendix-46da295c6787-definition-and-mechanism}

Phase 7 derives \((RS)^{-T}=RS^{-1}\) for orthogonal \(R\) and nonzero diagonal scales \(S\), explains why zero scale invalidates inverse-based normal transport, and distinguishes shading-normal direction from determinant-signed cross-derived area. \href{../02-Mathematics-for-Visual-Computing/main.pdf\#nameddest=01-foundations-i-ch019-normals-undo-scale-before-following-rotation}{TRS normal derivation (Mathematics for Visual Computing, Ch. 3)}.

\end{allmightroledefinition}

A covector is a linear scalar-valued measurement rule. Surface normals represent rules annihilating tangent movement. Under an invertible active linear transform, preserved measurements require inverse-transpose normal transformation. A metric identifies vector and covector representations.

For tangent \((1,1,0)\) and normal \((1,-1,0)\), doubling horizontal scale requires normal \((1/2,-1,0)\) to preserve perpendicularity.

\subsection{Prerequisites}\label{appendix-46da295c6787-prerequisites}

\begin{itemize}
\tightlist
\item
  \hyperref[appendix-ce0341f7ea0d]{Basis coordinates and dimension} --- An ordered basis fixes coordinate meaning and the space whose geometry is measured.
\item
  \href{../06-GPU-Architecture-and-Real-Time-Rendering/main.pdf\#nameddest=appendix-8bc4e8fc547c}{Dot products and orthogonality (GPU Architecture and Real-Time Rendering, Reference notes: Dot products and orthogonality)} --- Inner-product measurements supply the perpendicularity and alignment constraints.
\end{itemize}

\subsection{Used by}\label{appendix-46da295c6787-used-by}

Later curriculum sources can deepen this same node; no future coverage is marked complete.

\subsection{Recall and next depth}\label{appendix-46da295c6787-recall-and-next-depth}

Derive inverse transpose, distinguish active/passive changes, and handle singular transformations explicitly. Full production proofs, target-specific integration and learner mastery require further evidence.

\subsection{Source and navigation}\label{appendix-46da295c6787-source-and-navigation}

\href{../02-Mathematics-for-Visual-Computing/main.pdf\#nameddest=01-foundations-i-ch018-a-normal-is-a-measuring-rule-before-it-is-an-arrow}{Teaching narrative (Mathematics for Visual Computing, Ch. 2)} · \hyperref[appendix-fd090d982a15-phase-6---parts-501-600]{Coverage and source gap} · \hyperref[appendix-a103e7846d21-phase-6---source-analysis-and-verification]{Evidence and limits} · \hyperref[appendix-4ad47bdd269a]{Dependency Map}
